% Generated draft: requires mathematical and source review.
\documentclass[11pt]{article}
\usepackage[margin=25mm]{geometry}
\usepackage{fontspec}
\setmainfont{Latin Modern Roman}
\usepackage{amsmath,amssymb,mathtools,unicode-math}
\setmathfont{Latin Modern Math}
\usepackage{enumitem,needspace,xurl,hyperref}
\setlength{\parindent}{0pt}
\setlength{\parskip}{5pt}
\setlength{\emergencystretch}{3em}
\newcommand{\N}{\mathbb{N}}
\newcommand{\Z}{\mathbb{Z}}
\newcommand{\Q}{\mathbb{Q}}
\newcommand{\R}{\mathbb{R}}
\newcommand{\C}{\mathbb{C}}
\newcommand{\F}{\mathbb{F}}
\newcommand{\A}{\mathbb{A}}
\newcommand{\PP}{\mathbb{P}}
\newcommand{\E}{\mathbb{E}}
\newcommand{\eps}{\varepsilon}
\newcommand{\dd}{\,\mathrm{d}}
\begin{document}
\section*{UnsolvedMath ID: 30001891}
\textbf{Source problem code:} catalogue entry
\subsection{Definitions and mathematical statement}
Let \(d\geq 0\) be an integer. A convex \(d\)-polytope is the convex hull of finitely many points in \(\R^d\), with affine hull \(\R^d\). Its combinatorial type is the isomorphism class of its face lattice, including the empty face and the whole polytope. Two such polytopes are projectively equivalent if a real projective transformation in \(\operatorname{PGL}_{d+1}(\R)\), defined throughout the first polytope in the chosen affine chart, maps it onto the second. A polytope \(P\) is \emph{projectively unique} if every convex \(d\)-polytope with a face lattice isomorphic to that of \(P\) is projectively equivalent to \(P\).

\textbf{Finiteness question, using the supplied corrected formulation.} For each fixed \(d\), is the set
\[
\mathcal U_d=\left\{[\mathcal F(P)] : P\subset\R^d\text{ is a projectively unique convex }d\text{-polytope}\right\}
\]
finite, where \([\mathcal F(P)]\) denotes the isomorphism class of the face lattice of \(P\)? The zero-dimensional case is included by convention; the supplied dated assessment singles out \(d\geq4\) as the intended unresolved range. This assessment is not independently verified here.

\textbf{Realization-space question and unresolved conventions.} To describe one possible realization-space convention, let \(K\) be the face lattice of a realizable \(d\)-polytope whose \(n\) vertices have fixed labels. Define
\[
\mathcal R_{\mathrm{vert}}(K)=\left\{(v_1,\ldots,v_n)\in(\R^d)^n : \operatorname{conv}\{v_1,\ldots,v_n\}\text{ has dimension }d\text{ and labeled face lattice }K\right\}.
\]
Here the \(v_i\) must be distinct and must be exactly the vertices of the convex hull. This is an unquotiented vertex-coordinate realization space, whose dimension can be understood as semialgebraic dimension. Alternatively, realizations may be identified under affine or projective transformations, producing different realization-space conventions.

The supplied broader question asks how realization-space dimension grows with polytope size. It does not specify the realization-space convention, whether size means the number of vertices, facets, all faces, or another quantity, or whether the desired growth concerns upper bounds, lower bounds, extrema, or individual families. It also does not explicitly fix the dimension regime for this broader question. Consequently, this part supplies a research question rather than a uniquely specified quantified asymptotic statement. The displayed vertex-coordinate space is an explanatory convention, not an assertion that the source selected it.

The original wording combines combinatorial types with equivalence under projective transformations without imposing projective uniqueness. The supplied corrected record restores that hypothesis; the finiteness question above follows that correction rather than treating the original and corrected wordings as equivalent.
\subsection{Short English statement}
For each fixed dimension, ask whether only finitely many combinatorial types of projectively unique convex polytopes exist. The accompanying question about growth of realization-space dimension requires clarification of the realization-space convention, the measure of polytope size, and the intended bounds. The supplied open-status assessment requires independent review.
\subsection{Sources}
\begin{itemize}
\item Supplied local UnsolvedMath registry record, permanent ID \(30001891\), titled \emph{Projectively Unique Polytopes and Their Realization Spaces}; exact source problem code \texttt{catalogue entry}. The record belongs to the Oberwolfach Reports Open Problems collection and supplies proposed year 2011, with no named proposer.
\item Supplied source citation: \emph{Discrete Geometry (2012)}, Oberwolfach Reports, pp.~2459--2548, DOI \texttt{10.4171/owr/2011/44}. \url{https://doi.org/10.4171/owr/2011/44}. The supplied literature-source label is ``OWR 44 (2012).''
\item Additional supplied literature source, labeled ``unique polyhedral representations'': \url{https://arxiv.org/abs/math/0401005}. No authors or further bibliographic details are supplied in the record.
\item Literature assessment dated exactly \texttt{2026-08-21}: the supplied curation labels the problem open, describes the projectively unique finiteness question for \(d\geq4\) as open, and describes low dimensions as classified. This is supplied dated curation, not an independently verified literature assessment. The background explicitly requests independent verification of time-sensitive status and cited claims.
\item The record marks its statement \texttt{corrected\_verified} and reports that the official report was checked to restore the projective-uniqueness condition and realization-space question. These are supplied verification claims; no source visit or independent verification is asserted in this draft.
\end{itemize}
\end{document}
